% ============================================================================
%  Appendix E -- Answers to the questions (generated by tools/gen_appendices.py
%  from the slide decks, where they are the speaker notes of the quiz slides)
% ============================================================================
\chapter{Answers to the Questions}
\label{app:answers}

Each topical and tutorial chapter ends with multiple-choice questions in the
style of the midterm test (section \emph{Check your understanding}). The
answers follow: the letter, the text of the right option, and a line of
explanation. Answer first, and read afterwards: a question whose answer has
been read teaches nothing.

\section*{Chapter 1: Setting Up}
\begin{description}[leftmargin=3.2em,labelwidth=2.6em,itemsep=3pt]
\answer{1.1}{B}{\emph{ESP-IDF is not active in this terminal.} Activation lasts as long as the terminal: run export.sh, or open the ESP-IDF PowerShell.}
\answer{1.2}{C}{\emph{PRG32-QT.} PRG32-QT interprets the cartridge and has a built-in debugger; a compiler is needed only to build cartridges.}
\answer{1.3}{C}{\emph{it calls the runtime through the same table of functions on every host.} A portable cartridge depends only on the ABI, the table of functions of the runtime, which every host provides.}
\end{description}

\section*{Chapter 2: First Contact}
\begin{description}[leftmargin=3.2em,labelwidth=2.6em,itemsep=3pt]
\answer{2.1}{B}{\emph{The address of the counter.} la loads the address of the label; lw then reads the value stored at that address into t1.}
\answer{2.2}{C}{\emph{ra, because call overwrites it.} call writes the return address in ra. The callee's own calls would overwrite it, so a non-leaf function saves ra on the stack and restores it before ret.}
\answer{2.3}{D}{\emph{44.} 300 = 256 + 44 = 1 0010 1100 in binary. AND with 1111 1111 keeps the low eight bits: 0010 1100 = 44.}
\end{description}

\section*{Chapter 3: History of Computing}
\begin{description}[leftmargin=3.2em,labelwidth=2.6em,itemsep=3pt]
\answer{3.1}{B}{\emph{Program and data are stored in the same memory.} The stored-program concept: instructions are data in memory, fetched and executed one after another.}
\answer{3.2}{C}{\emph{Its instruction set is an open standard, free to implement.} RISC-V is an open standard instruction set architecture; anyone may design and sell a RISC-V processor without an architecture licence. It was not the first RISC (Berkeley RISC, MIPS, ARM came earlier).}
\answer{3.3}{C}{\emph{Ada Lovelace.} Ada Lovelace, for her 1843 notes on the Analytical Engine.}
\end{description}

\section*{Chapter 4: Digital Information}
\begin{description}[leftmargin=3.2em,labelwidth=2.6em,itemsep=3pt]
\answer{4.1}{D}{\emph{256.} A byte is 8 bits: $2^{8}$ = 256 patterns, for example the naturals from 0 to 255.}
\answer{4.2}{B}{\emph{magenta.} 0xF81F = 11111 000000 11111: red and blue at maximum, green at zero: magenta.}
\answer{4.3}{C}{\emph{5.} Four characters plus the terminating zero byte.}
\answer{4.4}{C}{\emph{0xCD.} RISC-V is little-endian: the least significant byte, 0xCD, goes at the lowest address.}
\end{description}

\section*{Chapter 5: Integers}
\begin{description}[leftmargin=3.2em,labelwidth=2.6em,itemsep=3pt]
\answer{5.1}{B}{\emph{-4.} Complement and add one: 0000 0011 + 1 = 0000 0100 = 4, so the pattern is -4. (252 is its unsigned reading.)}
\answer{5.2}{B}{\emph{100 + 28.} The range is [-128, 127]. 100 + 27 = 127 fits; 100 + 28 = 128 does not: two positives give a negative pattern (1000 0000 = -128).}
\answer{5.3}{C}{\emph{t2 = 1, t3 = 0.} Signed: -1 $<$ 1 is true, so t2 = 1. Unsigned: 4,294,967,295 $<$ 1 is false, so t3 = 0.}
\answer{5.4}{A}{\emph{0.101.} 0.625 $\times$ 2 = 1.25 (1); 0.25 $\times$ 2 = 0.5 (0); 0.5 $\times$ 2 = 1.0 (1): 0.101 = 1/2 + 1/8.}
\end{description}

\section*{Chapter 6: Real Numbers}
\begin{description}[leftmargin=3.2em,labelwidth=2.6em,itemsep=3pt]
\answer{6.1}{D}{\emph{Excess-k.} In excess-k, X = x + k: a larger number always has a larger pattern. This is why it encodes floating-point exponents, which must be compared quickly.}
\answer{6.2}{B}{\emph{1.5.} 0x180 = 384; 384 / 256 = 1.5. Integer part 0x01, fraction 0x80 = 128/256 = 0.5.}
\answer{6.3}{B}{\emph{3.0.} 0100 0000 0100 0000 ...: S = 0, E = 1000 0000 = 128, F = 100...: $1.1_2\times 2^{128-127}=1.5\times 2=3.0$.}
\answer{6.4}{C}{\emph{infinity.} All ones in the exponent: infinity if the fraction is zero, NaN otherwise.}
\end{description}

\section*{Chapter 7: Verm1, Part 1}
\begin{description}[leftmargin=3.2em,labelwidth=2.6em,itemsep=3pt]
\answer{7.1}{B}{\emph{4.} 64/256 = 1/4 of a cell per frame: a whole cell every 4 frames.}
\answer{7.2}{C}{\emph{(80, 32).} x = 10 $\times$ 8 = 80; y = 3 $\times$ 8 + 8 = 32.}
\answer{7.3}{B}{\emph{draw calls other functions, which overwrite ra.} A call writes the return address in ra. A function that makes calls must keep its own ra safe on the stack; a leaf function need not.}
\end{description}

\section*{Chapter 8: Boolean Algebra: Operators and Properties}
\begin{description}[leftmargin=3.2em,labelwidth=2.6em,itemsep=3pt]
\answer{8.1}{C}{\emph{x = 1, y = 1.} NAND is the complement of AND: it is 0 only when both inputs are 1.}
\answer{8.2}{B}{\emph{$\nott{a}+b$.} $\nott{a\cdot\nott{b}}=\nott{a}+\nott{\nott{b}}=\nott{a}+b$, using involution on the second term.}
\answer{8.3}{C}{\emph{0x0000000A.} 0x5A = 0101 1010; AND with 0000 1111 keeps the low nibble: 1010 = 0xA.}
\answer{8.4}{B}{\emph{ori t0, t0, 0x20.} OR with a mask: x + 1 = 1 on bit 5, x + 0 = x elsewhere. xori would toggle it; addi would carry into bit 6 if the bit was already 1.}
\end{description}

\section*{Chapter 9: Boolean Algebra: Functions and Minimisation}
\begin{description}[leftmargin=3.2em,labelwidth=2.6em,itemsep=3pt]
\answer{9.1}{D}{\emph{256.} $2^{2^{3}}=2^{8}=256$.}
\answer{9.2}{B}{\emph{$A\,\nott{B}\,C$.} Row 5 is 101: $A=1$, $B=0$, $C=1$, so the minterm is $A\,\nott{B}\,C$.}
\answer{9.3}{B}{\emph{two variables.} A group of $2^{k}$ cells eliminates k variables: four cells, two variables.}
\answer{9.4}{C}{\emph{at least two inputs are 1.} $C_{out}$ is the majority function: $AB+C_{in}(A\oplus B)=AB+AC_{in}+BC_{in}$. An odd number of 1s describes the sum bit S instead.}
\end{description}

\section*{Chapter 10: Electronics}
\begin{description}[leftmargin=3.2em,labelwidth=2.6em,itemsep=3pt]
\answer{10.1}{B}{\emph{260 $\Omega$.} R = (3.3 - 2) / 0.005 = 260 $\Omega$.}
\answer{10.2}{B}{\emph{the pMOS conducts and the output is 1.} A low gate closes the pMOS, which joins the output to the supply; the nMOS is open.}
\answer{10.3}{C}{\emph{6.} A NAND (4) followed by an inverter (2).}
\answer{10.4}{B}{\emph{To make the pin read a defined 1 when the button is released.} Without it the released pin would float; the pull-up ties it to 3.3 V, and the pressed switch overrides it to ground.}
\end{description}

\section*{Chapter 11: Combinational Machines}
\begin{description}[leftmargin=3.2em,labelwidth=2.6em,itemsep=3pt]
\answer{11.1}{B}{\emph{5.} $\lceil\log_2 26\rceil$ = 5, since $2^{4}$ = 16 $<$ 26 $\le$ 32 = $2^{5}$. Six words remain unassigned: the code is incomplete.}
\answer{11.2}{C}{\emph{8 outputs, one at 1.} $2^{3}$ = 8 outputs, and exactly one is asserted: the decoded (one-hot) representation of the input word.}
\answer{11.3}{B}{\emph{3.} n selection lines choose among $2^{n}$ inputs: 3 lines for 8 inputs.}
\answer{11.4}{C}{\emph{does not detect the error.} Two alterations leave the parity unchanged. A single parity bit detects only an odd number of errors.}
\end{description}

\section*{Chapter 12: Verm1, Part 2}
\begin{description}[leftmargin=3.2em,labelwidth=2.6em,itemsep=3pt]
\answer{12.1}{C}{\emph{0x0a.} 0x08 OR 0x02 = 0x0a: one bit per button.}
\answer{12.2}{B}{\emph{DOWN (1).} 3 XOR 1 = 2: DOWN is the opposite of UP.}
\answer{12.3}{B}{\emph{each entry of the table is 4 bytes.} The entries are words: entry k is at offset 4k.}
\answer{12.4}{C}{\emph{t1 and t4 are equal.} XOR is 1 where bits differ: no difference, all zeros.}
\end{description}

\section*{Chapter 13: Sequential Machines}
\begin{description}[leftmargin=3.2em,labelwidth=2.6em,itemsep=3pt]
\answer{13.1}{B}{\emph{the present state only.} Moore: w maps S to U. In a Mealy machine the output depends on the state and on the input.}
\answer{13.2}{C}{\emph{Q stays at 1.} S = R = 0 is the hold condition: the feedback keeps the stored value.}
\answer{13.3}{C}{\emph{The flip-flop changes only on a clock edge.} The latch is transparent for the whole active level of the clock; the flip-flop samples D only on the edge.}
\answer{13.4}{B}{\emph{4.} $\lceil\log_2 10\rceil$ = 4 state variables, with the minimum-length binary code of Chapter~11. Ten would be needed with a one-hot code.}
\end{description}

\section*{Chapter 14: Registers and Buses}
\begin{description}[leftmargin=3.2em,labelwidth=2.6em,itemsep=3pt]
\answer{14.1}{B}{\emph{So that only the enabled source drives the lines.} Several sources share the same lines; all but one must be electrically disconnected (high impedance), otherwise two outputs would fight over the wire.}
\answer{14.2}{C}{\emph{32.} One bit per pulse: n pulses for n bits. A parallel transfer needs one pulse and n wires.}
\answer{14.3}{C}{\emph{8,192.} $2^{10}$ = 1,024 words of 8 bits: 8,192 bits, that is 1 KiB.}
\end{description}

\section*{Chapter 15: Memory}
\begin{description}[leftmargin=3.2em,labelwidth=2.6em,itemsep=3pt]
\answer{15.1}{C}{\emph{64 KiB.} $2^{16}$ = 65,536 locations of one byte: 64 KiB, the address space of a Commodore 64.}
\answer{15.2}{D}{\emph{0x01.} Little-endian: the least significant byte, 0x01, is at the lowest address.}
\answer{15.3}{A}{\emph{Because its capacitors lose their charge.} Each bit is a charge in a capacitor that leaks away; a circuit inside the memory rewrites every cell periodically.}
\answer{15.4}{B}{\emph{With sw to the address of a device register.} RISC-V has no I/O instructions: devices are memory-mapped and reached with ordinary loads and stores.}
\end{description}

\section*{Chapter 16: Permanent Memory}
\begin{description}[leftmargin=3.2em,labelwidth=2.6em,itemsep=3pt]
\answer{16.1}{D}{\emph{Flash.} Flash is the block-erasable evolution of EEPROM. EPROM needs ultraviolet light; ROM and PROM cannot be erased.}
\answer{16.2}{B}{\emph{The datum is fetched from main memory and copied into the cache, with its neighbours.} The access is simply slower: the block containing the address is loaded into the cache, betting that nearby addresses will be needed soon.}
\answer{16.3}{B}{\emph{Latency.} Latency is the wait for the sector to arrive under the head; seek time depends on the movement of the arm.}
\answer{16.4}{C}{\emph{a transition between pit and land.} A change of reflectivity, in either direction, is a 1; no change is a 0.}
\end{description}

\section*{Chapter 17: Verm1, Part 3}
\begin{description}[leftmargin=3.2em,labelwidth=2.6em,itemsep=3pt]
\answer{17.1}{C}{\emph{254, 255, 0, 1.} From the tail, head\_i - length + 1 = -2, that is 254 after AND 255, up to the head: 254, 255, 0, 1.}
\answer{17.2}{B}{\emph{byte 127.} 3 $\times$ 40 + 7 = 127.}
\answer{17.3}{B}{\emph{the game starts once.} Only the frame in which A goes down gives s1 = 0x10: one edge, one transition.}
\answer{17.4}{B}{\emph{so that the head can enter the cell the tail leaves.} In the same step the tail moves away: that cell is free for the head.}
\end{description}

\section*{Chapter 18: The Processor}
\begin{description}[leftmargin=3.2em,labelwidth=2.6em,itemsep=3pt]
\answer{18.1}{B}{\emph{the status register with the flags.} RISC-V has no condition codes: branches compare registers directly.}
\answer{18.2}{C}{\emph{only loads and stores.} It is a load-store architecture: lw, lh, lb, sw, sh, sb and their unsigned forms.}
\answer{18.3}{B}{\emph{t0 was pushed on the stack.} The pointer moves down by one word and the value is stored on the new top: a push.}
\answer{18.4}{B}{\emph{accumulator := accumulator + [y].} The accumulator is the implicit operand and the destination.}
\end{description}

\section*{Chapter 19: Instructions and Interrupts}
\begin{description}[leftmargin=3.2em,labelwidth=2.6em,itemsep=3pt]
\answer{19.1}{B}{\emph{the seven least significant.} Bits 6 to 0, in every format: the processor always knows where to look first.}
\answer{19.2}{C}{\emph{4.} Addresses count bytes and an instruction is four bytes long.}
\answer{19.3}{C}{\emph{mepc.} mepc, written by the hardware at the trap and used by mret. mtvec holds the address of the handler.}
\answer{19.4}{C}{\emph{an illegal instruction is fetched.} It is caused by the instruction being executed, synchronously; the others come from outside, at any moment.}
\end{description}

\section*{Chapter 20: Buses and the Video Card}
\begin{description}[leftmargin=3.2em,labelwidth=2.6em,itemsep=3pt]
\answer{20.1}{C}{\emph{921,600 bytes.} 307,200 pixels $\times$ 3 bytes = 921,600 bytes.}
\answer{20.2}{B}{\emph{100 MB/s.} 2 bytes $\times$ 50 million transfers per second = 100 MB/s.}
\answer{20.3}{C}{\emph{the position of its colour in a palette.} The byte is an index; the actual colour is looked up in a table of 256 entries, the palette.}
\answer{20.4}{B}{\emph{Using graphics processors for non-graphical computation.} General-purpose computing on graphics processing units.}
\end{description}

\section*{Chapter 21: Input and Output}
\begin{description}[leftmargin=3.2em,labelwidth=2.6em,itemsep=3pt]
\answer{21.1}{B}{\emph{It turns a row and column combination into a key code.} From a decoded position in the key matrix to a compact n-bit code, as the encoder of Chapter~11.}
\answer{21.2}{B}{\emph{three cells with red, green and blue filters.} Three sub-pixels, each a liquid crystal cell behind a colour filter; a TFT in each stores its state.}
\answer{21.3}{B}{\emph{Its pixels are driven one at a time and must hold their state in between.} Persistence: without a transistor per pixel, the cells respond slowly and keep their state badly while the others are scanned.}
\answer{21.4}{C}{\emph{polling.} The program asks the device periodically. With interrupts the device would call the program.}
\end{description}

\section*{Chapter 22: Verm1, Part 4}
\begin{description}[leftmargin=3.2em,labelwidth=2.6em,itemsep=3pt]
\answer{22.1}{B}{\emph{32 30 35.} The characters '2', '0', '5': 0x32, 0x30, 0x35. The units are stored at offset 2.}
\answer{22.2}{B}{\emph{8.} One step every 4 frames is 64/256: 32 + 4 $\times$ score = 64, so score = 8.}
\answer{22.3}{B}{\emph{5.} la is two instructions, auipc and addi; then lw, beqz and ret.}
\answer{22.4}{B}{\emph{it calls no function and uses only a and t registers.} A leaf function that touches only caller-saved registers has nothing to save.}
\end{description}

\section*{Chapter 23: The Cartridge Store}
\begin{description}[leftmargin=3.2em,labelwidth=2.6em,itemsep=3pt]
\answer{23.1}{B}{\emph{publish 1.0.1.} A published version is treated as permanent: a correction is a new version with the same id.}
\answer{23.2}{B}{\emph{the metadata goes inside one cartridge, which has one architecture.} The list of files belongs to the bundle; a single cartridge states its own architecture, given on the command line.}
\answer{23.3}{C}{\emph{in an environment variable of your shell.} A token is a credential: it never goes into a file that is shared or versioned.}
\end{description}

\section*{Chapter 24: Beyond One Core}
\begin{description}[leftmargin=3.2em,labelwidth=2.6em,itemsep=3pt]
\answer{24.1}{B}{\emph{5.} With n unlimited, S = 1 / (1 - p) = 1 / 0.2 = 5.}
\answer{24.2}{C}{\emph{false dependences between uses of the same register name.} Two unrelated values that share a name get different physical registers and can be computed at the same time.}
\answer{24.3}{B}{\emph{30.} 256 bits are 32 bytes per turn: 960 / 32 = 30. The code does not change.}
\answer{24.4}{B}{\emph{OR.} The output is 1 whenever at least one input is 1: OR.}
\end{description}
